\section{Procedure}
Sometimes called \emph{Method} or \emph{Experimental}.

Report the operating procedure that was actually followed, concisely but completely, so that the experiment can be reproduced. A schematic drawing of the apparatus helps when the assembly is not obvious.

Comments and observations are welcome when they are relevant, for instance details not mentioned in the given procedure but noticed during the experiment.

\info[inline]{Example: Add Devarda's alloy. \textit{Warning: after the addition the test tube becomes very hot; work near the bench and keep a beaker at hand to rest the tube in.}}

A flowchart that summarises the steps is also useful, both for clarity and for studying (see the flowchart example at the end of this document).
